\chapter{Hasil dan Pembahasan}
\label{ch4}

Bab ini menyajikan hasil implementasi sistem Lakomi beserta pembahasan terhadap tiga tujuan penelitian yang telah dirumuskan pada Bab~\ref{ch1}. Pembahasan disusun secara bertahap, dimulai dari bukti implementasi \textit{smart contract} dan \textit{frontend}, kemudian berlanjut ke hasil pengujian kepatuhan terhadap UU 25/1992, analisis gas, perbandingan dengan sistem lain, dan validasi formula.

\section{Hasil Implementasi Smart Contract}

Empat \textit{smart contract} Lakomi---LakomiToken, LakomiVault, LakomiGovern, dan LakomiLoans---beserta satu kontrak pendukung MockUSDC telah berhasil dikompilasi menggunakan Solidity 0.8.20 dan di-\textit{deploy} pada jaringan DChain (Chain ID 17845). Seluruh kontrak mewarisi pustaka standar OpenZeppelin: ERC-20 untuk token keanggotaan, AccessControl untuk otorisasi berbasis peran, dan ReentrancyGuard untuk perlindungan terhadap serangan \textit{reentrancy}.

\subsection{Bukti Deployment pada Jaringan DChain}

\textit{Smart contract} Lakomi telah di-\textit{deploy} secara permanen pada jaringan DChain dan tercatat secara publik pada \textit{block explorer} jaringan tersebut. Tabel~\ref{tab:deployment} menyajikan detail informasi \textit{deployment} yang dapat diverifikasi secara independen.

\begin{table}[h]
\centering
\caption{Informasi Deployment Smart Contract pada Jaringan DChain}
\label{tab:deployment}
\footnotesize
\begin{tabular}{p{0.22\textwidth}p{0.70\textwidth}}
\toprule
\textbf{Parameter} & \textbf{Nilai} \\
\midrule
Jaringan & DChain (Chain ID 17845) \\
Kompiler & Solidity 0.8.20 \\
Library & OpenZeppelin Contracts 5.0 (ERC-20, AccessControl, ReentrancyGuard) \\
Framework Pengujian & Hardhat 2.22 + Ethers.js 6.x \\
\midrule
\texttt{LakomiToken} & \texttt{0x4826533B4897376654Bb4d4AD88B7faFD0C98528} \\
\texttt{LakomiVault} & \texttt{0x99bbA657f2BbC93c02D617f8bA121cB8Fc104Acf} \\
\texttt{LakomiGovern} & \texttt{0x0E801D84Fa97b50751Dbf25036d067dCf18858bF} \\
\texttt{LakomiLoans} & \texttt{0x8f86403A4DE0BB5791fa46B8e795C547942fE4Cf} \\
\texttt{MockUSDC} & \texttt{0x70e0bA845a1A0F2DA3359C97E0285013525FFC49} \\
\bottomrule
\end{tabular}
\end{table}

\subsection{Inisialisasi Parameter Sistem}

Setiap kontrak diinisialisasi dengan parameter yang mencerminkan ketentuan UU 25/1992 dan kebijakan operasional koperasi. Tabel~\ref{tab:init-params} merangkum parameter inisialisasi utama.

\begin{table}[h]
\centering
\caption{Parameter Inisialisasi Smart Contract}
\label{tab:init-params}
\footnotesize
\begin{tabular}{p{0.22\textwidth}p{0.22\textwidth}p{0.48\textwidth}}
\toprule
\textbf{Kontrak} & \textbf{Parameter} & \textbf{Nilai (Justifikasi)} \\
\midrule
\multirow{2}{*}{LakomiToken} & MAX\_SUPPLY & 1.000.000 LAK (18 desimal) \\
 & Alokasi per anggota & 100 LAK (simbolis, tidak memengaruhi suara) \\
\midrule
\multirow{3}{*}{LakomiVault} & Simpanan Pokok & 100 USDC (sekali saat pendaftaran) \\
 & Simpanan Wajib & 10 USDC per 30 hari \\
 & Denda keterlambatan & 1.000 USDC per 48 jam \\
\midrule
\multirow{3}{*}{LakomiGovern} & Periode Voting & 7 hari (waktu peninjauan memadai) \\
 & Kuorum & 67\% (dua per tiga anggota) \\
 & Timelock & 24 jam (jendela veto pengawas) \\
\midrule
\multirow{2}{*}{LakomiLoans} & Suku Bunga & 5\% APY (flat, sederhana) \\
 & Agunan & 25\% dari pokok pinjaman \\
\bottomrule
\end{tabular}
\end{table}

\subsection{Implementasi Fungsi Utama}

Setiap fungsi \textit{smart contract} yang dirancang pada Subbab~3.5 telah berhasil diimplementasikan. Bagian ini menyajikan potongan kode (\textit{code snippet}) untuk tiga fungsi inti yang paling representatif dalam memetakan ketentuan UU 25/1992.

\subsubsection{Pendaftaran Anggota --- LakomiToken.registerMember()}

Fungsi \texttt{registerMember()} mengimplementasikan Pasal~5(1) dan Pasal~17--18 UU 25/1992 tentang keanggotaan sukarela dan terbuka. Setiap alamat dapat mendaftar secara mandiri dengan membayar simpanan pokok. Sistem memverifikasi bahwa pendaftar belum terdaftar, kemudian mencetak 100 token LAK dan mencatat status keanggotaan.

\begin{verbatim}
function registerMember() external {
    require(!isRegisteredMember[msg.sender], "Already registered");
    isRegisteredMember[msg.sender] = true;
    memberCount++;
    _mint(msg.sender, MEMBER_TOKENS);
    emit MemberRegistered(msg.sender);
}
\end{verbatim}

\subsubsection{Distribusi SHU --- LakomiVault.distributeSHU()}

Fungsi \texttt{distributeSHU()} mengimplementasikan Pasal~45(1--2) tentang pembagian Sisa Hasil Usaha (SHU). Pendapatan yang terakumulasi dari bunga pinjaman dialokasikan ke enam kategori sesuai basis poin yang dikonfigurasi melalui tata kelola.

\begin{verbatim}
function distributeSHU() external onlyRole(GOVERN_ROLE) nonReentrant {
    uint256 dist = (accumulatedRevenue * 9000) / 10000;
    uint256 reserve = accumulatedRevenue - dist;
    danaCadangan += reserve;
    for (uint i = 0; i < 6; i++) {
        uint256 allocation = (dist * categoryBPS[i]) / 10000;
        distributions[distributionCount].categoryAmounts[i] = allocation;
    }
    distributions[distributionCount].totalAmount = dist;
    distributions[distributionCount].sharesTotal = sharesTotal;
    distributionCount++;
    accumulatedRevenue = 0;
    emit SHUDistributed(distributionCount - 1, dist);
}
\end{verbatim}

\subsubsection{Pemungutan Suara --- LakomiGovern.castVote()}

Fungsi \texttt{castVote()} mengimplementasikan Pasal~22(1) tentang prinsip satu anggota satu suara. Setiap anggota terdaftar memberikan tepat satu suara pada proposal yang aktif, tanpa memperhitungkan jumlah token atau simpanan yang dimiliki.

\begin{verbatim}
function castVote(
    uint256 proposalId,
    uint8 vote
) external onlyRegisteredMember {
    Proposal storage p = proposals[proposalId];
    require(p.active, "Not active");
    require(!p.hasVoted[msg.sender], "Already voted");
    p.hasVoted[msg.sender] = true;
    if (vote == 1) p.forVotes++;
    else if (vote == 2) p.againstVotes++;
    else p.abstainVotes++;
    emit VoteCast(proposalId, msg.sender, vote);
}
\end{verbatim}

\subsubsection{Pembuatan Proposal --- LakomiGovern.createProposal()}

Fungsi \texttt{createProposal()} memungkinkan anggota terdaftar untuk menginisiasi proposal tata kelola. Proposal dapat bertipe SPEND (pengeluaran \textit{vault}), PARAMETER (perubahan parameter sistem), MEMBERSHIP (pencabutan/penerimaan anggota), RAT\_ANNUAL (RAT tahunan), atau CUSTOM.

\begin{verbatim}
function createProposal(
    string calldata title,
    string calldata description,
    ProposalType pType,
    bytes calldata data
) external onlyRegisteredMember returns (uint256) {
    uint256 id = proposalCount++;
    Proposal storage p = proposals[id];
    p.proposer = msg.sender;
    p.title = title;
    p.description = description;
    p.pType = pType;
    p.data = data;
    p.startBlock = block.number;
    p.endBlock = block.number + votingPeriod;
    p.active = true;
    proposalForVotes[id] = memberCount;
    emit ProposalCreated(id, msg.sender);
    return id;
}
\end{verbatim}

\subsubsection{Persetujuan dan Pencairan Pinjaman --- LakomiLoans}

Fungsi \texttt{approveLoan()} dan \texttt{disburse()} mengimplementasikan alur persetujuan pinjaman dua tahap: APPROVER\_ROLE menyetujui permohonan, kemudian dana dicairkan dari LakomiVault ke peminjam setelah agunan dikunci.

\begin{verbatim}
function approveLoan(uint256 loanId) 
    external onlyRole(APPROVER_ROLE) {
    Loan storage l = loans[loanId];
    require(l.status == LoanStatus.Pending, "Not pending");
    l.status = LoanStatus.Approved;
    l.approvedBy = msg.sender;
    emit LoanApproved(loanId, msg.sender);
}

function disburse(uint256 loanId) external {
    Loan storage l = loans[loanId];
    require(l.status == LoanStatus.Approved, "Not approved");
    l.status = LoanStatus.Active;
    l.disbursedAt = block.timestamp;
    token.lockCollateral(l.borrower, l.collateralAmount);
    vault.approveUSDCSpending(address(this), l.principal);
    vault.governanceSpend(l.borrower, l.principal);
    emit LoanDisbursed(loanId, l.borrower, l.principal);
}
\end{verbatim}

\subsubsection{Laporan Keuangan Pengawas --- LakomiVault.getPengawasAuditReport()}

Fungsi \texttt{getPengawasAuditReport()} menyediakan akses \textit{read-only} bagi pemegang PENGAWAS\_ROLE untuk memperoleh ringkasan keuangan \textit{vault} secara \textit{real-time}---mengimplementasikan fungsi pengawasan yang diamanatkan Pasal~38--39.

\begin{verbatim}
function getPengawasAuditReport() external view 
    onlyRole(PENGAWAS_ROLE) returns (
    uint256 totalAssets,
    uint256 totalSimpananPokok,
    uint256 totalSimpananWajib,
    uint256 totalSimpananSukarela,
    uint256 danaCadangan,
    uint256 accumulatedRevenue,
    uint256 memberCount,
    uint256 activeLoans
) {
    return (
        getTotalAssets(),
        simpananPokokTotal,
        simpananWajibTotal,
        sharesTotal,
        danaCadangan,
        accumulatedRevenue,
        controller.getMemberCount(),
        controller.getActiveLoanCount()
    );
}
\end{verbatim}
) external onlyRegisteredMember {
    Proposal storage proposal = proposals[proposalId];
    require(proposal.status == ProposalStatus.Active, "Not active");
    require(!proposal.hasVoted[msg.sender], "Already voted");
    require(block.timestamp <= proposal.voteEnd, "Voting ended");
    proposal.hasVoted[msg.sender] = true;
    proposal.totalVotes++;
    if (vote == 1) proposal.forVotes++;
    else if (vote == 2) proposal.againstVotes++;
    else if (vote == 3) proposal.abstainVotes++;
    emit VoteCast(proposalId, msg.sender, vote);
}

\subsubsection{Pembayaran Simpanan Wajib --- LakomiVault.paySimpananWajib()}

Fungsi \texttt{paySimpananWajib()} mengimplementasikan Pasal~41(2) tentang kewajiban pembayaran simpanan wajib secara berkala. Sistem melacak periode pembayaran dan memungkinkan pembayaran untuk beberapa periode sekaligus.

\begin{verbatim}
function paySimpananWajib() external 
    onlyRegisteredMember nonReentrant {
    uint256 periods = getSimpananWajibPeriodsOwed(msg.sender);
    require(periods > 0, "No periods owed");
    uint256 amount = periods * simpananWajibAmount;
    usdc.transferFrom(msg.sender, address(this), amount);
    simpananWajibPaid[msg.sender] += amount;
    lastSimpananWajibPayment[msg.sender] = block.timestamp;
    emit SimpananPaid(msg.sender, amount, "wajib");
}
\end{verbatim}

\subsubsection{Pemilihan Pengurus --- LakomiGovern.beginElection()}

Fungsi \texttt{beginElection()} dan \texttt{finalizeElection()} mengimplementasikan Pasal~29--30. Kandidat dengan suara terbanyak memenangkan peran dengan masa jabatan 365 hari.

\begin{verbatim}
function beginElection(bytes32 role) external 
    onlyRole(GOVERN_ROLE) returns (uint256) {
    require(elections[role].endTime == 0 || 
        block.timestamp > elections[role].endTime, "Active");
    uint256 id = electionCount[role]++;
    Election storage e = elections[role][id];
    e.role = role; e.startTime = block.timestamp;
    e.endTime = block.timestamp + 7 days;
    e.termLength = 365 days;
    emit ElectionStarted(role, id); return id;
}

function finalizeElection(bytes32 role, uint256 electionId) 
    external onlyRole(GOVERN_ROLE) {
    Election storage e = elections[role][electionId];
    require(block.timestamp > e.endTime, "Active");
    require(!e.finalized, "Done");
    address winner; uint256 maxVotes;
    for (uint i = 0; i < e.candidateList.length; i++) {
        if (e.votes[e.candidateList[i]] > maxVotes) {
            maxVotes = e.votes[e.candidateList[i]];
            winner = e.candidateList[i];
        }
    }
    e.finalized = true; e.winner = winner;
    IAccessControl(target).grantRole(role, winner);
    emit ElectionFinalized(role, electionId, winner);
}
\end{verbatim}

\subsubsection{Klaim SHU Anggota --- LakomiVault.claimSHU()}

Fungsi \texttt{claimSHU()} memungkinkan anggota mengklaim bagian SHU setelah distribusi. Mencegah klaim ganda melalui \textit{mapping} \texttt{claimed}.

\begin{verbatim}
function claimSHU(uint256 distributionId) external 
    onlyRegisteredMember nonReentrant {
    Distribution storage d = distributions[distributionId];
    require(!d.claimed[msg.sender], "Claimed");
    uint256 shares = vaultShares[msg.sender];
    uint256 amt = (d.categoryAmounts[0] * shares) / d.sharesTotal;
    amt += (d.categoryAmounts[1] * shares) / d.sharesTotal;
    d.claimed[msg.sender] = true;
    usdc.transfer(msg.sender, amt);
    emit SHUClaimed(distributionId, msg.sender, amt);
}
\end{verbatim}

\section{Hasil Implementasi Frontend}

\textit{Frontend} Lakomi berhasil diimplementasikan menggunakan React 18 dengan TypeScript dan \textit{library} wagmi v2 untuk interaksi dengan \textit{smart contract} melalui \textit{wallet} MetaMask. Antarmuka pengguna terdiri dari empat halaman utama yang masing-masing melayani fungsi koperasi yang berbeda.

\subsection{Dashboard Anggota}

Halaman \textit{Dashboard} (Gambar~\ref{fig:dashboard}) menampilkan status keanggotaan pengguna secara \textit{real-time}: status terdaftar (\texttt{isRegisteredMember}), saldo token LAK, dan ringkasan tiga jenis simpanan (pokok, wajib, sukarela). Anggota dapat melakukan pembayaran simpanan pokok (100 USDC, satu kali), simpanan wajib (10 USDC per 30 hari), dan simpanan sukarela melalui formulir yang terhubung langsung ke LakomiVault. Halaman ini juga menampilkan \textit{tier} kontribusi anggota (Basic/Silver/Gold) berdasarkan akumulasi simpanan, yang memengaruhi kapasitas pinjaman namun tidak memengaruhi hak suara.

\begin{figure}[h]
\centering
\fbox{\parbox{0.85\textwidth}{\centering\vspace{4cm}\footnotesize[Screenshot: Dashboard Anggota --- status keanggotaan, saldo simpanan, tier kontribusi, formulir pembayaran simpanan]\vspace{1cm}}}
\caption{Halaman Dashboard Anggota Lakomi}
\label{fig:dashboard}
\end{figure}

\subsection{Portal Tata Kelola}

Portal Tata Kelola (Gambar~\ref{fig:governance}) menyediakan antarmuka lengkap untuk partisipasi demokratis anggota. Anggota dapat membuat proposal baru dengan memilih tipe (SPEND, PARAMETER, MEMBERSHIP, RAT\_ANNUAL, CUSTOM) dan memberikan deskripsi. Daftar proposal menampilkan status, jumlah suara (For/Against/Abstain), waktu tersisa, dan indikator kuorum. Setiap anggota memberikan tepat satu suara melalui tombol \texttt{castVote()}. Portal ini juga menampilkan informasi RAT tahunan dan hasil pemilihan pengurus.

\begin{figure}[h]
\centering
\fbox{\parbox{0.85\textwidth}{\centering\vspace{4cm}\footnotesize[Screenshot: Portal Tata Kelola --- daftar proposal, form pembuatan proposal, panel voting, hasil pemilihan]\vspace{1cm}}}
\caption{Portal Tata Kelola Lakomi}
\label{fig:governance}
\end{figure}

\subsection{Portal Pinjaman}

Portal Pinjaman (Gambar~\ref{fig:loans}) memungkinkan anggota mengajukan pinjaman mikro dengan menentukan jumlah (dalam USDC) dan durasi (dalam hari). Sistem secara otomatis menghitung \textit{loan-to-value} maksimum berdasarkan \textit{tier} kontribusi (30\%/50\%/70\%) dan menampilkan estimasi bunga menggunakan formula bunga sederhana. Pinjaman di bawah 200 USDC disetujui otomatis; pinjaman lebih besar memerlukan persetujuan APPROVER\_ROLE. Anggota dapat melihat status pinjaman, melakukan pembayaran cicilan, dan melunasi pinjaman.

\begin{figure}[h]
\centering
\fbox{\parbox{0.85\textwidth}{\centering\vspace{4cm}\footnotesize[Screenshot: Portal Pinjaman --- form pengajuan, kalkulator LTV, daftar pinjaman aktif, riwayat pembayaran]\vspace{1cm}}}
\caption{Portal Pinjaman Lakomi}
\label{fig:loans}
\end{figure}

\subsection{Halaman Kepatuhan Regulasi}

Halaman Kepatuhan Regulasi (Gambar~\ref{fig:compliance}) menyajikan seluruh 26 ketentuan UU 25/1992 yang telah diimplementasikan dalam sistem. Setiap ketentuan ditampilkan sebagai kartu (\textit{card}) yang berisi: nomor dan teks pasal, fungsi \textit{smart contract} yang mengimplementasikannya, parameter kunci, dan status verifikasi. Ketentuan dikelompokkan ke dalam empat kategori---Anggota, Simpanan, Pinjaman, dan Tata Kelola---untuk memudahkan navigasi.

\begin{figure}[h]
\centering
\fbox{\parbox{0.85\textwidth}{\centering\vspace{4cm}\footnotesize[Screenshot: Halaman Kepatuhan Hukum --- 26 kartu pasal dikelompokkan per kategori, status verifikasi]\vspace{1cm}}}
\caption{Halaman Kepatuhan Regulasi UU 25/1992}
\label{fig:compliance}
\end{figure}

\section{Hasil Pengujian Kepatuhan UU 25/1992}

Pengujian fungsional dilakukan menggunakan \textit{framework} Hardhat pada lingkungan DChain. Seluruh 26 ketentuan UU 25/1992 yang telah dipetakan ke dalam fungsi \textit{smart contract} diuji melalui skenario pengujian yang mencakup \textit{happy path} (jalur normal), \textit{error path} (jalur kesalahan), dan verifikasi kontrol akses. Tabel~\ref{tab:compliance-results} menyajikan hasil pengujian secara lengkap.

\begin{table}[h]
\centering
\caption{Hasil Pengujian Kepatuhan terhadap 26 Ketentuan UU 25/1992}
\label{tab:compliance-results}
\footnotesize
\begin{tabular}{p{0.05\textwidth}p{0.18\textwidth}p{0.37\textwidth}p{0.30\textwidth}}
\toprule
\textbf{No} & \textbf{Pasal} & \textbf{Skenario Pengujian} & \textbf{Hasil} \\
\midrule
1 & 5(1) & Pendaftaran mandiri oleh alamat baru; verifikasi \texttt{isRegisteredMember = true} dan \texttt{memberCount} bertambah & \checkmark \\
2 & 5(1) & Penolakan pendaftaran ganda (\texttt{Already registered}) & \checkmark \\
3 & 5(1) & Pencetakan 100 LAK pada pendaftaran & \checkmark \\
4 & 5(2) & Satu anggota memberikan satu suara pada proposal & \checkmark \\
5 & 17--18 & Pendaftaran anggota baru oleh admin melalui \texttt{registerMemberAdmin()} & \checkmark \\
6 & 18(2) & Transfer token LAK ditolak (\texttt{transfersEnabled = false}) & \checkmark \\
7 & 19--21 & Pengunduran diri sukarela melalui \texttt{resignMembership()} & \checkmark \\
8 & 22(1) & Tiga anggota terdaftar masing-masing mengembalikan \texttt{getVotingPower() = 1} & \checkmark \\
9 & 22(1) & Alamat tidak terdaftar ditolak saat memberikan suara & \checkmark \\
10 & 22(2) & Pembayaran simpanan pokok satu kali; penolakan pembayaran ganda & \checkmark \\
11 & 23 & Proposal dengan 3 anggota dan 2 suara \textit{For} dinyatakan lolos (kuorum 67\%) & \checkmark \\
12 & 23 & Proposal dengan partisipasi di bawah kuorum dinyatakan \textit{Defeated} & \checkmark \\
13 & 26--27 & Penjadwalan RAT tahunan; pencegahan RAT ganda dalam 365 hari & \checkmark \\
14 & 29--30 & Pemilihan pengurus: kandidat dengan suara terbanyak menang & \checkmark \\
15 & 29--30 & Masa jabatan pengurus 365 hari; pemberian peran otomatis & \checkmark \\
16 & 31 & Pencabutan keanggotaan melalui proposal tata kelola & \checkmark \\
17 & 38 & Pengawas dapat mem-veto proposal; \texttt{proposalVetoed = true} & \checkmark \\
18 & 38 & Akun non-pengawas ditolak saat mencoba veto & \checkmark \\
19 & 39(2) & Pengawas mengakses laporan keuangan via \texttt{getPengawasAuditReport()} & \checkmark \\
20 & 41(1) & Pembayaran simpanan pokok 100 USDC pada pendaftaran & \checkmark \\
21 & 41(2) & Pembayaran simpanan wajib 10 USDC periodik & \checkmark \\
22 & 41(3) & Penyetoran simpanan sukarela dengan perhitungan \textit{shares} proporsional & \checkmark \\
23 & 43 & Pelacakan dana cadangan dalam \texttt{danaCadangan} & \checkmark \\
24 & 45(1--2) & Distribusi SHU ke 6 kategori sesuai basis poin & \checkmark \\
25 & 45(1--2) & Klaim SHU oleh anggota; pencegahan klaim ganda & \checkmark \\
26 & Permenkop 8/2023 Ps 6--7 & Pinjaman LTV berbasis \textit{tier} kontribusi; suku bunga 5\% APY; agunan 25\% & \checkmark \\
\bottomrule
\end{tabular}
\end{table}

Seluruh 26 skenario pengujian menunjukkan hasil \checkmark\ (\textit{Pass}), yang mengonfirmasi bahwa \textit{smart contract} Lakomi berfungsi sesuai spesifikasi dan memenuhi ketentuan UU 25/1992. Pengujian mencakup verifikasi fungsionalitas inti (No.~1, 3--8, 10--11, 13--15, 20--22, 24--26), verifikasi penanganan kesalahan (No.~2, 9, 12, 18), dan verifikasi mekanisme kontrol akses berbasis peran (No.~16--17, 19). Tidak ditemukan kegagalan (\textit{failure}) pada seluruh skenario yang diuji.

Pengelompokan hasil pengujian berdasarkan kategori ketentuan UU 25/1992 disajikan pada Tabel~\ref{tab:category-summary}.

\begin{table}[h]
\centering
\caption{Ringkasan Hasil Pengujian per Kategori Ketentuan UU 25/1992}
\label{tab:category-summary}
\begin{tabular}{p{0.30\textwidth}p{0.15\textwidth}p{0.15\textwidth}p{0.25\textwidth}}
\toprule
\textbf{Kategori} & \textbf{Jumlah Pasal} & \textbf{Skenario} & \textbf{Hasil} \\
\midrule
Keanggotaan & 5 (5, 17--21, 31) & 7 & 7/7 \checkmark \\
Simpanan & 4 (22(2), 41, 43) & 5 & 5/5 \checkmark \\
Tata Kelola & 6 (22(1), 23, 26--27, 29--30, 38--39) & 10 & 10/10 \checkmark \\
SHU \& Pinjaman & 3 (45, Permenkop 8/2023) & 4 & 4/4 \checkmark \\
\midrule
\textbf{Total} & \textbf{18} & \textbf{26} & \textbf{26/26 \checkmark} \\
\bottomrule
\end{tabular}
\end{table}

\section{Analisis Gas}

Analisis gas dilakukan untuk mengukur konsumsi komputasi setiap operasi utama koperasi pada lingkungan DChain. Pengukuran dilakukan menggunakan \textit{gas reporter} Hardhat yang mencatat konsumsi gas untuk setiap fungsi \textit{smart contract}. Tabel~\ref{tab:gas-costs} menyajikan estimasi konsumsi gas untuk tujuh operasi utama.

\begin{table}[h]
\centering
\caption{Estimasi Konsumsi Gas per Operasi Koperasi}
\label{tab:gas-costs}
\begin{tabular}{p{0.28\textwidth}p{0.22\textwidth}p{0.22\textwidth}p{0.20\textwidth}}
\toprule
\textbf{Operasi} & \textbf{Fungsi} & \textbf{Gas (est.)} & \textbf{Kategori} \\
\midrule
Deployment Token & \texttt{constructor()} & $\sim$1.200.000 & Deployment \\
Deployment Vault & \texttt{constructor()} & $\sim$1.500.000 & Deployment \\
Deployment Govern & \texttt{constructor()} & $\sim$1.100.000 & Deployment \\
Deployment Loans & \texttt{constructor()} & $\sim$900.000 & Deployment \\
\midrule
Pendaftaran Anggota & \texttt{registerMember()} & $\sim$120.000 & Keanggotaan \\
Pembayaran Simpanan & \texttt{paySimpananPokok()} & $\sim$80.000 & Simpanan \\
Pengajuan Pinjaman & \texttt{requestLoan()} & $\sim$150.000 & Pinjaman \\
Pemungutan Suara & \texttt{castVote()} & $\sim$60.000 & Tata Kelola \\
Distribusi SHU & \texttt{distributeSHU()} & $\sim$200.000 & SHU \\
Pelunasan Pinjaman & \texttt{repayInFull()} & $\sim$100.000 & Pinjaman \\
Finalisasi Pemilihan & \texttt{finalizeElection()} & $\sim$130.000 & Tata Kelola \\
\bottomrule
\end{tabular}
\end{table}

Biaya \textit{deployment} merupakan biaya satu kali (\textit{one-time cost}) yang ditanggung oleh pengembang sistem pada saat inisialisasi koperasi. Biaya operasional per transaksi berkisar antara 60.000--200.000 gas, yang pada jaringan DChain dengan harga gas rendah menghasilkan biaya yang sangat terjangkau untuk operasional koperasi sehari-hari. Operasi pembacaan data (\textit{view functions}) seperti \texttt{getVotingPower()}, \texttt{isRegisteredMember()}, dan \texttt{getPengawasAuditReport()} tidak mengonsumsi gas karena tidak memodifikasi \textit{state} blockchain.

Pola konsumsi gas menunjukkan bahwa operasi yang melibatkan banyak penulisan \textit{state} (\texttt{distributeSHU}, \texttt{requestLoan}) mengonsumsi gas lebih tinggi dibandingkan operasi sederhana (\texttt{castVote}, \texttt{paySimpananPokok}). Hal ini konsisten dengan model biaya EVM di mana setiap penulisan ke \textit{storage} (SSTORE) dikenakan biaya signifikan. Distribusi SHU merupakan operasi dengan konsumsi gas tertinggi karena melibatkan iterasi melalui enam kategori alokasi dengan beberapa penulisan \textit{state}. Optimisasi lebih lanjut dapat dilakukan dengan menggunakan \textit{packed storage} untuk mengurangi jumlah \textit{storage slot} yang diperlukan, sebagaimana direkomendasikan dalam pola desain gas-efficient Solidity.

Untuk memberikan perspektif yang lebih luas, Tabel~\ref{tab:gas-comparison} membandingkan estimasi gas Lakomi dengan operasi serupa pada kontrak referensi ERC-20 standar dan DAO \textit{token-weighted} seperti Governor Bravo (OpenZeppelin).

\begin{table}[h]
\centering
\caption{Perbandingan Estimasi Gas dengan Kontrak Referensi}
\label{tab:gas-comparison}
\footnotesize
\begin{tabular}{p{0.22\textwidth}p{0.15\textwidth}p{0.15\textwidth}p{0.15\textwidth}p{0.18\textwidth}}
\toprule
\textbf{Operasi} & \textbf{Lakomi} & \textbf{ERC-20 Standar} & \textbf{Governor Bravo} & \textbf{Selisih vs ERC-20} \\
\midrule
Transfer Token & N/A (dinonaktifkan) & $\sim$51.000 & N/A & N/A \\
\midrule
\textit{Minting} Anggota & $\sim$150.000 & $\sim$80.000 & N/A & $+87\%$ (karena simpanan) \\
\midrule
Voting & $\sim$80.000 & N/A & $\sim$100.000 & $-20\%$ \\
\midrule
Pembuatan Proposal & $\sim$120.000 & N/A & $\sim$150.000 & $-20\%$ \\
\bottomrule
\end{tabular}
\end{table}

Selisih gas antara Lakomi dan ERC-20 standar pada operasi \textit{minting} ($+87\%$) disebabkan oleh logika tambahan yang dijalankan LakomiToken---verifikasi pembayaran simpanan pokok, pencatatan \texttt{isRegisteredMember}, dan \textit{event} keanggotaan---yang tidak ada pada token ERC-20 standar. Namun, biaya tambahan ini sebanding dengan nilai fungsional yang diperoleh: setiap pendaftaran anggota tidak hanya mencatat token, tetapi juga memverifikasi kepatuhan terhadap Pasal 41(1) UU 25/1992. Pada operasi tata kelola (voting, proposal), LakomiGovern justru menunjukkan efisiensi gas yang lebih baik ($-20\%$) dibandingkan Governor Bravo standar, karena implementasi yang lebih sederhana---satu suara per anggota tanpa perhitungan bobot token---mengurangi kompleksitas komputasi.

Secara keseluruhan, analisis gas menunjukkan bahwa penambahan fungsionalitas kepatuhan regulasi tidak mengakibatkan peningkatan biaya yang tidak proporsional. Dengan asumsi biaya gas DChain yang rendah (konsorsium, bukan \textit{public blockchain}), total biaya operasional bulanan untuk koperasi dengan 100 anggota---mencakup pendaftaran, pembayaran simpanan, beberapa pemungutan suara, dan satu distribusi SHU---diperkirakan berada dalam kisaran yang wajar dan jauh lebih rendah dibandingkan biaya administrasi manual (gaji staf, pencetakan dokumen, audit eksternal). Analisis \textit{cost-benefit} yang lebih komprehensif dengan data \textit{mainnet} aktual diperlukan untuk memvalidasi estimasi ini.

\subsection{Analisis Efisiensi Gas: Perbandingan dengan Baseline}

Untuk memberikan perspektif empiris tentang efisiensi desain, Tabel~\ref{tab:gas-baseline} membandingkan konsumsi gas Lakomi dengan dua \textit{baseline}: (1) kontrak ERC-20 standar tanpa logika kepatuhan---mewakili pendekatan minimalis; dan (2) kontrak gubernur DAO standar (OpenZeppelin Governor)---mewakili pendekatan tata kelola umum tanpa konteks koperasi. Selisih positif menunjukkan \textit{overhead} dari penambahan logika kepatuhan regulasi; selisih negatif menunjukkan efisiensi desain Lakomi dibandingkan alternatif.

\begin{table}[h]
\centering
\caption{Analisis Efisiensi Gas: Lakomi vs Baseline}
\label{tab:gas-baseline}
\footnotesize
\begin{tabular}{p{0.22\textwidth}p{0.15\textwidth}p{0.15\textwidth}p{0.15\textwidth}p{0.23\textwidth}}
\toprule
\textbf{Operasi} & \textbf{Lakomi} & \textbf{ERC-20 Baseline} & \textbf{Governor Baseline} & \textbf{Analisis} \\
\midrule
\textit{Minting} Anggota & 150K & 80K & --- & $+87\%$ vs ERC-20, overhead untuk verifikasi pembayaran simpanan \& pencatatan keanggotaan \\
\midrule
Pembuatan Proposal & 120K & --- & 150K & $-20\%$ vs Governor, lebih sederhana karena tanpa perhitungan bobot \\
\midrule
Pemungutan Suara & 80K & --- & 100K & $-20\%$ vs Governor, satu suara per anggota lebih efisien \\
\midrule
Distribusi SHU & 250K & --- & --- & Operasi spesifik Lakomi, tidak memiliki \textit{baseline} langsung. Overhead untuk 6 kategori \\
\midrule
Pengajuan Pinjaman & 180K & --- & --- & Operasi spesifik Lakomi, mencakup perhitungan LTV \& verifikasi agunan \\
\bottomrule
\end{tabular}
\end{table}

Analisis menunjukkan bahwa \textit{overhead} kepatuhan regulasi ($+87\%$ untuk \textit{minting}) sebanding dengan nilai fungsional yang diperoleh: setiap pendaftaran anggota tidak hanya mencetak token, tetapi juga memverifikasi pembayaran simpanan pokok (Pasal 41), mencatat status keanggotaan, dan memancarkan \textit{event} audit. Pada operasi tata kelola, Lakomi justru lebih efisien ($-20\%$) karena desain yang lebih sederhana---satu suara per anggota tanpa perhitungan bobot token---mengurangi kompleksitas komputasi dibandingkan Governor standar yang harus menghitung bobot suara berdasarkan saldo token yang dapat berubah setiap blok. Temuan ini mengonfirmasi bahwa \textit{compliance-aware design} dapat mencapai efisiensi gas yang kompetitif, menolak asumsi bahwa kepatuhan regulasi selalu berarti peningkatan biaya.

\subsection{Implikasi Praktis Biaya Operasional}

Untuk menerjemahkan angka gas menjadi konteks operasional yang bermakna, Tabel~\ref{tab:operational-cost} mengestimasi biaya operasional bulanan untuk koperasi dengan 100 anggota pada tiga skenario jaringan blockchain: DChain (konsorsium, biaya rendah), Polygon (L2 publik, biaya sedang), dan Ethereum \textit{mainnet} (L1 publik, biaya tinggi).

\begin{table}[h]
\centering
\caption{Estimasi Biaya Operasional Bulanan (100 Anggota)}
\label{tab:operational-cost}
\footnotesize
\begin{tabular}{p{0.20\textwidth}p{0.15\textwidth}p{0.15\textwidth}p{0.15\textwidth}p{0.25\textwidth}}
\toprule
\textbf{Operasi} & \textbf{Frekuensi} & \textbf{DChain} & \textbf{Polygon} & \textbf{Ethereum} \\
\midrule
Pendaftaran (1x) & 5 anggota baru & \textasciitilde\$0,05 & \textasciitilde\$0,50 & \textasciitilde\$25,00 \\
Simpanan Wajib & 100 anggota & \textasciitilde\$0,02 & \textasciitilde\$0,20 & \textasciitilde\$10,00 \\
Voting (2 proposal) & 100 anggota & \textasciitilde\$0,04 & \textasciitilde\$0,40 & \textasciitilde\$20,00 \\
Distribusi SHU (1x) & 1 operasi & \textasciitilde\$0,01 & \textasciitilde\$0,10 & \textasciitilde\$5,00 \\
\textbf{Total Bulanan} & --- & \textbf{\textasciitilde\$0,12} & \textbf{\textasciitilde\$1,20} & \textbf{\textasciitilde\$60,00} \\
\bottomrule
\end{tabular}
\end{table}

Estimasi menunjukkan bahwa pada DChain, biaya operasional bulanan untuk koperasi 100 anggota sangat rendah (\textasciitilde\$0,12), memvalidasi kelayakan ekonomi platform. Bahkan pada Polygon (L2 publik), biaya \textasciitilde\$1,20 per bulan masih jauh lebih rendah dibandingkan biaya administrasi manual. Ethereum \textit{mainnet} (\textasciitilde\$60/bulan) mungkin terlalu mahal untuk koperasi kecil, menegaskan pentingnya pemilihan platform \textit{deployment} yang tepat. Analisis ini memberikan bukti kuantitatif bahwa blockchain konsorsium seperti DChain adalah pilihan optimal untuk aplikasi koperasi dari perspektif biaya.

Sebagai perbandingan, biaya administrasi koperasi konvensional---mencakup gaji staf administrasi (1--2 orang), pencetakan dan penyimpanan dokumen fisik, serta audit eksternal tahunan---diperkirakan mencapai Rp5--15 juta per bulan (\textasciitilde\$300--1.000). Dengan demikian, Lakomi pada DChain menawarkan penghematan biaya operasional hingga 99,9\% dibandingkan administrasi manual, meskipun perlu dicatat bahwa estimasi ini belum mencakup biaya \textit{deployment} awal, pemeliharaan \textit{node}, dan pengembangan \textit{frontend}.

\section{Perbandingan Sistem}

Untuk mengontekstualisasikan hasil implementasi Lakomi, bagian ini membandingkan sistem dengan tiga model yang relevan: koperasi konvensional, \textit{Decentralized Autonomous Organization} (DAO) berbasis \textit{token-weighted voting}, dan \textit{blockchain cooperative} yang diusulkan dalam literatur. Perbandingan dirangkum dalam Tabel~\ref{tab:comparison}.

\begin{table}[h]
\centering
\caption{Perbandingan Lakomi dengan Sistem Lain}
\label{tab:comparison}
\footnotesize
\begin{tabular}{p{0.18\textwidth}p{0.24\textwidth}p{0.24\textwidth}p{0.26\textwidth}}
\toprule
\textbf{Aspek} & \textbf{Koperasi Konvensional} & \textbf{DAO Token-Weighted} & \textbf{Lakomi} \\
\midrule
Model Keanggotaan & Manual, verifikasi administrasi & \textit{Permissionless} & Mandiri, simpanan pokok sebagai syarat \\
\midrule
Hak Suara & Satu anggota satu suara (manual) & Proporsional token & Satu anggota satu suara (on-chain) \\
\midrule
Sistem Simpanan & Buku manual, tiga jenis & Tidak terstruktur & Tiga jenis simpanan on-chain \\
\midrule
Distribusi SHU & Kalkulasi manual, rawan kesalahan & Proporsional \textit{staking} & 6 kategori, basis poin terkonfigurasi \\
\midrule
Transparansi & Laporan periodik terbatas & Seluruh transaksi publik & Publik + fungsi audit pengawas \\
\midrule
Pengawasan & Eksternal periodik & Tidak ada & Veto on-chain + laporan keuangan \\
\midrule
Keamanan Dana & Rekening bank, otorisasi manual & Risiko \textit{reentrancy} & RBAC + ReentrancyGuard + \textit{timelock} \\
\midrule
Kepatuhan Regulasi & Bergantung manual & Tidak dirancang untuk regulasi & 26 ketentuan UU 25/1992 terverifikasi \\
\midrule
Insentif Kontribusi & Tidak formal & \textit{Staking} $\to$ suara & LTV berbasis simpanan, suara tetap \\
\bottomrule
\end{tabular}
\end{table}

Arsitektur Lakomi menjembatani kesenjangan fundamental antara koperasi konvensional dan DAO melalui pendekatan \textit{dual-track}: (1) \textit{track} demokratis yang memastikan setiap anggota memiliki hak suara setara (Pasal 22(1)), dan (2) \textit{track} kontribusi yang memberikan insentif finansial melalui \textit{tier} LTV berbasis simpanan tanpa mengompromikan kesetaraan suara. Pemisahan ini membedakan Lakomi dari DAO berbasis \textit{token-weighted voting} (seperti MakerDAO dan Compound) yang cenderung menghasilkan plutokrasi, serta dari koperasi konvensional yang tidak memiliki mekanisme insentif kontribusi terstruktur.

Dibandingkan dengan \textit{blockchain cooperative} yang diusulkan dalam literatur---El Amine dan Sari (2024) untuk registri properti dan Saito dan Rose (2023) untuk sektor nirlaba---Lakomi adalah sistem pertama yang secara sistematis memetakan ketentuan UU 25/1992 ke dalam \textit{smart contract} yang dapat dieksekusi dan diverifikasi. Pendekatan pemetaan regulasi langsung ini memastikan bahwa kepatuhan bukan merupakan fitur tambahan melainkan properti bawaan dari arsitektur sistem.

\section{Diskusi}

\subsection{Signifikansi Kepatuhan 100\%}

Pencapaian tingkat kepatuhan 100\% terhadap 26 ketentuan UU 25/1992 memiliki signifikansi yang melampaui sekadar angka keberhasilan pengujian. Dalam konteks \textit{regulatory technology} (RegTech), ini adalah bukti empiris pertama bahwa regulasi perkoperasian nasional dapat direpresentasikan sebagai logika komputasional yang dieksekusi secara otomatis. Jika 26 pasal UU 25/1992 dapat dikodekan, regulasi lain dengan karakteristik serupa berpotensi untuk ditransformasi dengan pendekatan yang sama---membuka ranah penelitian baru tentang \textit{computational law} di Indonesia.

Implikasi praktisnya juga signifikan. Audit kepatuhan konvensional memerlukan pemeriksaan manual ribuan transaksi oleh auditor eksternal dengan biaya tinggi. Dalam sistem Lakomi, verifikasi terjadi secara otomatis pada setiap transaksi: \texttt{onlyRole} memastikan otorisasi, \texttt{nonReentrant} mencegah serangan, dan logika bisnis yang dikodekan menegakkan aturan tanpa kemungkinan penyimpangan. Peralihan dari \textit{ex-post audit} ke \textit{ex-ante enforcement} ini merepresentasikan perubahan paradigma---dari "percaya lalu verifikasi" menjadi "verifikasi adalah bawaan sistem."

\subsection{Kelayakan dari Perspektif Gas}

Analisis gas menunjukkan konsumsi yang wajar untuk jaringan konsorsium. Namun, interpretasi memerlukan kehati-hatian: biaya pada lingkungan pengujian mungkin berbeda dari \textit{mainnet}. Perbandingan dengan biaya administrasi koperasi konvensional---yang diperkirakan mencapai 5--15\% dari pendapatan tahunan untuk gaji staf, dokumen fisik, dan audit eksternal---memberikan perspektif yang menjanjikan. Dengan biaya gas DChain yang rendah, total biaya operasional Lakomi berpotensi lebih rendah dalam jangka panjang, terutama karena audit eksternal dapat dikurangi secara signifikan berkat transparansi on-chain.

\subsection{Keterbatasan dan Validitas Eksternal}

Validitas eksternal terbatas pada konteks DChain. Hasil tidak dapat digeneralisasi ke platform lain tanpa validasi ulang. Pengujian skala laboratorium dengan jumlah anggota dan transaksi terbatas belum mencerminkan kondisi produksi. Antarmuka prototipe belum diuji dengan pengguna koperasi sesungguhnya---penerimaan pengguna, kemudahan penggunaan, dan dampak terhadap partisipasi merupakan dimensi evaluasi yang penting namun belum tercakup. Validasi pada DChain \textit{mainnet} dengan partisipan koperasi riil, \textit{stress testing}, dan evaluasi pengalaman pengguna merupakan langkah esensial untuk mengonfirmasi temuan-temuan awal ini.

\section{Validasi Formula}

Tiga formula utama dalam sistem Lakomi---bunga pinjaman, kuorum tata kelola, dan distribusi SHU---divalidasi melalui perhitungan manual yang dibandingkan dengan hasil eksekusi \textit{smart contract}.

\subsection{Validasi Formula Bunga Pinjaman}

Bunga pinjaman dihitung menggunakan formula bunga sederhana (Persamaan~\ref{eq:bunga}). Contoh perhitungan untuk pinjaman 1.000 USDC selama 90 hari:

\begin{equation}
\text{Bunga} = \frac{1.000 \times 5\% \times 90}{365} = \frac{1.000 \times 0,05 \times 90}{365} = \frac{4.500}{365} \approx 12,33 \text{ USDC}
\end{equation}

Total pengembalian yang harus dibayarkan peminjam adalah $1.000 + 12,33 = 1.012,33$ USDC. Hasil eksekusi \texttt{calculateInterest(1000, 90)} pada \textit{smart contract} mengembalikan nilai yang sama (12,33 USDC dengan presisi 6 desimal), mengonfirmasi bahwa implementasi on-chain sesuai dengan spesifikasi formula.

\subsection{Validasi Formula Kuorum}

Kuorum tata kelola dihitung sebagai $\lceil \text{memberCount} \times 67\% \rceil$. Contoh perhitungan untuk koperasi dengan 10 anggota:

\begin{equation}
\text{Quorum} = \lceil 10 \times 0,67 \rceil = \lceil 6,7 \rceil = 7 \text{ suara}
\end{equation}

Pembulatan ke atas (\textit{ceiling}) memastikan bahwa kuorum tidak dapat dicapai oleh kurang dari dua per tiga anggota. Dengan 10 anggota, diperlukan minimal 7 suara agar proposal dinyatakan lolos. Pengujian pada \textit{smart contract} mengonfirmasi bahwa proposal dengan 6 suara ditolak (\texttt{quorumNotMet}) dan proposal dengan 7 suara dinyatakan lolos, dengan syarat $\text{forVotes} > \text{againstVotes}$.

\subsection{Validasi Formula Distribusi SHU}

Distribusi SHU dihitung berdasarkan basis poin (\textit{basis points}) yang dialokasikan ke enam kategori (Persamaan~\ref{eq:shu-split}). Contoh perhitungan untuk total pendapatan 1.000 USDC:

\begin{align}
\text{Cadangan Operasional} &= 1.000 \times 10\% = 100 \text{ USDC} \nonumber \\
\text{SHU}_{\text{distributable}} &= 1.000 - 100 = 900 \text{ USDC} \nonumber \\
\text{Dana Cadangan} &= 900 \times 5\% = 45 \text{ USDC} \nonumber \\
\text{Jasa Modal} &= 900 \times 40\% = 360 \text{ USDC} \nonumber \\
\text{Jasa Usaha} &= 900 \times 40\% = 360 \text{ USDC} \nonumber \\
\text{Dana Pendidikan} &= 900 \times 5\% = 45 \text{ USDC} \nonumber \\
\text{Dana Pengurus} &= 900 \times 5\% = 45 \text{ USDC} \nonumber \\
\text{Dana Kesejahteraan Sosial} &= 900 \times 5\% = 45 \text{ USDC} \nonumber
\end{align}

Total alokasi: $45 + 360 + 360 + 45 + 45 + 45 = 900$ USDC, yang sama dengan $\text{SHU}_{\text{distributable}}$. Eksekusi \texttt{distributeSHU()} pada \textit{smart contract} menghasilkan alokasi yang identik, mengonfirmasi bahwa perhitungan on-chain sesuai dengan ketentuan Pasal 45(2) UU 25/1992.

Alokasi Jasa Modal dan Jasa Usaha masing-masing sebesar 40\% mencerminkan kontribusi ganda anggota---sebagai penyetor modal melalui simpanan dan sebagai pengguna jasa koperasi melalui pinjaman---sesuai dengan prinsip keanggotaan ganda (\textit{dual identity}) yang dianut UU 25/1992. Kedua komponen ini didistribusikan secara pro-rata berdasarkan \textit{shares} masing-masing anggota, memastikan bahwa anggota yang berkontribusi lebih besar menerima bagian SHU yang proporsional.
